\section{Related work}
\label{sec:related_work}

\subsection{Medical image segmentation}

U-Net \citep{ronneberger2015unet} and its descendants remain the dominant
architecture family for medical image segmentation. Encoder-decoder designs
with skip connections have been extended in numerous directions: attention
gates \citep{oktay2018attention}, dense connectivity \citep{li2018hdenseunet},
transformer-based encoders \citep{chen2021transunet}, and multi-scale feature
fusion \citep{zhou2018unetplusplus}. For retinal vessel segmentation
specifically, the DRIVE \citep{staal2004drive}, STARE \citep{hoover2000stare},
and CHASE\_DB1 \citep{fraz2012chase} benchmarks have driven steady
improvements in Dice and AUC, with recent methods exceeding 0.80 Dice on DRIVE.

Nearly all of this work optimizes aggregate segmentation metrics. A model
that achieves 0.78 Dice is considered better than one at 0.76, regardless of
whether the errors cluster on ambiguous images (where a clinician would want
to look more carefully) or scatter randomly. This paper takes a different
starting point: given a segmentation model of reasonable quality, how should
its outputs be used in a decision pipeline?

\subsection{Uncertainty estimation in deep learning}

Bayesian neural networks provide a principled framework for uncertainty but are
computationally prohibitive for modern architectures \citep{neal1996bayesian,
blundell2015weight}. Practical approximations fall into three families.

\paragraph{MC Dropout.} Gal and Ghahramani \citep{gal2016dropout} showed that
applying dropout at test time and averaging multiple stochastic forward passes
approximates variational inference. The variance (or mutual information) across
passes quantifies epistemic uncertainty. MC Dropout is widely used in medical
imaging \citep{nair2020exploring,jungo2018effect,roy2019bayesian} because it
requires no architectural changes and works with any dropout-equipped network.
Its weakness is that the quality of the uncertainty estimate depends on
where dropout layers are placed and on the number of forward passes $T$, with
diminishing returns beyond $T \approx 20$--30 in our experiments. The
uncertainty also tends to be diffuse: MC Dropout flags large regions as
somewhat uncertain rather than precisely localizing errors.

\paragraph{Deep ensembles.} Lakshminarayanan et al.\
\citep{lakshminarayanan2017simple} train $N$ models with different random
initializations and use prediction variance across ensemble members as
uncertainty. Ensembles consistently produce better-calibrated predictions than
single models \citep{ovadia2019can}, but training $N$ separate models multiplies
computational cost. In our experiments, a 5-member ensemble requires roughly
$4\times$ the inference time of a single MC Dropout model.

\paragraph{Test-Time Augmentation.} TTA generates multiple predictions by
applying geometric or photometric transformations to the input and aggregating
the (inverse-transformed) outputs \citep{wang2019aleatoric,
ayhan2018test,moshkov2020tta}. Disagreement across augmented predictions
indicates regions where the model's output is sensitive to input perturbations.
TTA captures a form of uncertainty that is complementary to MC Dropout:
where MC Dropout probes the model's weight posterior, TTA probes sensitivity
to input variations that should, in principle, not change the prediction. This
distinction is relevant because boundary regions and thin structures (common
error sources in vessel segmentation) tend to produce high TTA disagreement
even when MC Dropout uncertainty is moderate.

Mehrtash et al.\ \citep{mehrtash2020confidence} compared MC Dropout and TTA
for brain tumor segmentation and found that TTA-based uncertainty correlated
more strongly with segmentation error. Our work confirms this finding for
retinal vessels and extends it to the deferral setting, showing that TTA's
superior uncertainty quality translates directly to better deferral decisions.

\subsection{Calibration}

A model is well calibrated if its predicted probabilities match empirical
frequencies: among all pixels assigned probability 0.8, roughly 80\% should
be positive. Expected Calibration Error (ECE) \citep{naeini2015obtaining}
bins predictions by confidence and measures the average gap between predicted
and observed accuracy. Temperature scaling \citep{guo2017calibration} is the
standard post-hoc fix: a single scalar $T$ is learned on a validation set to
rescale logits before the sigmoid/softmax, shrinking or expanding the
confidence distribution.

The assumption behind calibration as a preprocessing step for decision-making
is that calibrated probabilities are more informative. This assumption is not
always correct. Calibration optimizes a global property (the probability
distribution matches reality on average) but deferral requires a local
property (high uncertainty should coincide with high error \emph{per pixel}).
A model can be perfectly calibrated yet have uncertainty that fails to
discriminate correct from incorrect predictions. We test this empirically
and find that temperature scaling barely moves ECE (from 0.038 to 0.038 for
MC Dropout, from 0.037 to 0.039 for TTA) and does not improve the tracked
deferral operating points in our main runs.

Laves et al.\ \citep{laves2020well} reached a similar conclusion for
classification, observing that calibration and selective prediction can be
at odds. Our work extends this observation to dense prediction (segmentation)
and to the specific setting of pixel-level deferral.

\subsection{Selective prediction and deferral}

Selective prediction, also called prediction with a reject option, has a long
history in classification \citep{chow1970optimum,el2010foundations}. The idea
is simple: a model outputs both a prediction and a confidence score, and
predictions below a confidence threshold are withheld. Geifman and El-Yaniv
\citep{geifman2017selective} formalized risk-coverage curves for deep
classifiers, showing that selective prediction can dramatically reduce error
rates at modest coverage costs.

Translating selective prediction to segmentation is harder because the unit
of abstention is unclear. Should the model abstain on individual pixels, on
patches, or on entire images? Pixel-level deferral is the most granular but
produces fragmented deferral maps that may be impractical for human review.
Image-level deferral is cleaner but discards good predictions along with
bad ones. We adopt pixel-level deferral but evaluate it at multiple
granularities, including per-image error reduction and coverage-averaged
metrics.

DeVries and Taylor \citep{devries2018learning} trained a separate confidence
branch alongside a classifier. Corbiere et al.\ \citep{corbiere2019addressing}
proposed learning a failure prediction module. In medical imaging, Jungo et al.\
\citep{jungo2020analyzing} studied how uncertainty thresholds affect the
quality of retained predictions in brain lesion segmentation, and
Mehrtash et al.\ \citep{mehrtash2020confidence} evaluated deferral using
MC Dropout for prostate segmentation.

Our confidence-aware deferral strategy differs from these approaches in that
it does not require training an auxiliary network. Instead, it combines the
existing uncertainty map with the model's own prediction confidence through a
simple multiplicative score, $s = u \cdot (1 - c)$, where $u$ is uncertainty
and $c$ is prediction confidence. This has a clear interpretation: a pixel
is deferred when it is both uncertain \emph{and} the model's prediction is
near the decision boundary. Uncertain pixels where the model is nonetheless
confident (far from $p = 0.5$) are retained. This design costs nothing
computationally and, as we show, substantially improves deferral efficiency.

The learning-to-defer framework of Madras et al.\ \citep{madras2018predict}
and Mozannar and Sontag \citep{mozannar2020consistent} formalizes deferral as
a joint optimization of model and expert, assuming known expert accuracy. Our
setting is simpler and arguably more practical: we assume no model of the
human expert and instead optimize deferral thresholds on a validation set to
maximize a task-relevant objective (error reduction at minimum coverage cost).
This makes our approach applicable to any existing segmentation pipeline
without retraining.
